\section{Evaluating use, retention, and adaptation}
\label{sec:capability_eval}
We evaluate agent research capability across three dimensions:

\paragraph{Method selection and task quality.}
Tasks present distinct quantitative problems requiring the agent to select tools, set
parameters, and interpret numerical outputs. Measuring whether chosen algorithms match input
data properties (stationarity, regime type, asset universe) alongside returns separates sound
reasoning from fortunate price paths across matched price, text, fixed-output, and autonomous controls.

\paragraph{Context retention and progressive routing.}
Holding the base model, input data, and tool interfaces fixed, we compare full-library
concatenation, fixed-window BM25 chunk retrieval, and two-stage progressive skill routing
(\texttt{RAG-Jev}) under strict point-in-time (\texttt{as\_of}) filtering across decision
quality, retention of parameter schemas and guard signatures, token footprint, and chunk fragmentation.

\paragraph{Experience and adaptation.}
Holding tool access and context budgets fixed, we evaluate sequential adaptation across
unaugmented models, frozen-memory baselines, and our multi-timescale memory module under
identical cleared returns, tracking post-loss method shifts, turnover costs, and regime
performance. All model outputs and portfolio weights are frozen prior to scoring.
